% Abstractseite
\newpage
\begin{center}\large\bfseries Abstract\end{center}

We apply a weighted Birkhoff average (WBA) to compute frequencies in a two and four dimensional setting.
We first provide a proof of concept by constructing artificial signals with known frequencies. %, and show how to overcome a problem in an oscillatory case.
The WBA appears to be superior to the numerical analysis of fundamental frequencies when used as an indication of chaotic dynamics, as well as in the context of trapped orbits in the two dimensional setting.
We demonstrate a problem, that arises in the strongly coupled case of a four dimensional standard map, and how it can be solved using a transformation obtained from diagonalizing the inertia tensor.
We then apply our solution to compute the frequency space of the map for a specific set of parameters, and show two types of orbits, for which the transformation only allows to compute a single frequency with WBA. 

\vspace*{11em}
\begin{center}\large\bfseries Zusammenfassung\end{center}

Wir verwenden ein gewichtetes Birkhoff-Mittel (WBA) zur Berechnung von Frequenzen in zwei und vier Dimensionen.
Wir liefern zunächst einen Proof of Concept, indem wir künstliche Signale mit bekannten Frequenzen konstruieren. %, und zeigen, wie ein Problem im oszillatorischen Fall gelöst werden kann.
Das WBA erweist sich gegenüber der numerischen Analyse fundamentaler Frequenzen in der Anwendung als Chaosindikator und der Frequenzanalyse gefangener Bahnen in zwei Dimensionen überlegen.
Wir demonstrieren ein Problem, das im stark gekoppelten Fall einer vierdimensionalen Standardabbildung auftritt, und wie es mit Hilfe einer Transformation gelöst werden kann, die sich aus der Diagonalisierung des Trägheitstensors ergibt.
Wir wenden unsere Lösung an, um den Frequenzraum der Abbildung für einen bestimmten Satz von Parametern zu berechnen, und zeigen zwei Arten von Bahnen, für welche die Transformation nur die Berechnung einer einzelnen Frequenz mit WBA ermöglicht. 
